\documentclass[12pt,a4paper]{article}
\usepackage[utf8]{inputenc}
\usepackage[T1]{fontenc}
\usepackage{enumitem}
\usepackage[margin=2.5cm]{geometry}
\usepackage{parskip}
\usepackage{amsmath}
\usepackage{amssymb}
\usepackage{xcolor}

\title{midterm --- Final Exam --- Answer Key}
\date{}

\begin{document}

\maketitle

{\large \textbf{\textcolor{red}{ANSWER KEY --- Not for Distribution}}}

\vspace{0.5em}



\noindent
\textbf{Total points:} 4
\hfill
\textbf{Number of questions:} 4

\vspace{1em}
\hrule
\vspace{1.5em}


\noindent
\textbf{Question 1 (1.00 pts).}~
What is a key difference between a classical bit and a qubit?


\begin{enumerate}[label=\textbf{\Alph*.}]

  \item A classical bit can be in a blend of states at once, while a qubit cannot.

  \item A qubit is like a light switch, always in exactly one of two positions.

  \item \textbf{\textcolor{teal}{[CORRECT]}}  A qubit is more like a coin spinning in the air, with an undecided outcome until measured.

  \item A classical bit is more like a coin spinning in the air, with an undecided outcome until measured.

\end{enumerate}



\textit{\small Explanation: A qubit is more like a coin spinning in the air, with an undecided outcome until measured, whereas a classical bit is like a light switch, always in exactly one of two positions.}



\vspace{0.8em}
\noindent\rule{\linewidth}{0.2pt}
\vspace{0.3em}


\noindent
\textbf{Question 2 (1.00 pts).}~
Quantum gates can destroy or create probability.


\textbf{Answer:} \textcolor{teal}{False}



\textit{\small Explanation: The context states that quantum gates do not destroy or create probability, they preserve it.}



\vspace{0.8em}
\noindent\rule{\linewidth}{0.2pt}
\vspace{0.3em}


\noindent
\textbf{Question 3 (1.00 pts).}~
What is the key difference between the correlation of entangled qubits and the correlation of matching socks?


\textbf{Model Answer:} \textcolor{teal}{The correlation of entangled qubits is created at the moment of measurement, whereas the correlation of matching socks is decided when they are packed.}



\textit{\small Explanation: Entangled qubits have no definite value until measured; Matching socks have their color decided when packed}



\vspace{0.8em}
\noindent\rule{\linewidth}{0.2pt}
\vspace{0.3em}


\noindent
\textbf{Question 4 (1.00 pts).}~
Explain the concept of superposition in quantum computing and how it differs from classical computing.


\textbf{Key Points:} \textcolor{teal}{I. Introduction to superposition, II. Comparison with classical computing, III. Implications of superposition in quantum computing}



\textit{\small Explanation: Definition of superposition (3 pts): Clearly defines superposition as a single, well-defined quantum state with no classical equivalent; Comparison with classical computing (4 pts): Accurately compares and contrasts superposition with classical computing concepts; Implications of superposition (3 pts): Effectively discusses the implications of superposition in quantum computing, such as its role in quantum gates and algorithms}





\end{document}
